\begin{lemma}[Exact localization of the transformed bad part, complex inputs]
Let $v, w \in L^2(\gamma_n;\mathbb{C})$, let
$g_1, g_2, u_1, u_2 \in L^2(\gamma_n;\mathbb{R})$ with
$u_1, u_2 \in \operatorname{Dom}(L^{1/2})$ in the sense of
\noderef{FracPowDomain}, and suppose that, with $B = L^{1/2}$ as in
\noderef{FracPow},
\[
  \operatorname{Re} v = g_1 + Bu_1, \qquad \operatorname{Im} v = g_2 + Bu_2
  \qquad\text{in } L^2(\gamma_n;\mathbb{R}),
\]
and that $w$ has these good parts as its own real and imaginary parts:
\[
  \operatorname{Re} w = g_1, \qquad \operatorname{Im} w = g_2 .
\]
Let $t \in \mathbb{R}$ and let $\Omega \subseteq \mathbb{R}^n$ be an arbitrary
subset such that $Du_1 = 0$ and $Du_2 = 0$ $\gamma_n$-a.e.\ on $\Omega^c$, with
$D$ as in \noderef{GaussD}; no measurability of $\Omega$ is required, and
$\gamma_n(\Omega)$ below is the outer measure, which agrees with the measure
when $\Omega$ is measurable.  Then, with $R_\gamma^{\mathbb{C}}$ as in
\noderef{ComplexRieszTransform}:
\begin{enumerate}
\item for $\gamma_n$-almost every $x$ with $x \notin \Omega$,
  $R_\gamma^{\mathbb{C}} v(x) = R_\gamma^{\mathbb{C}} w(x)$; and
\item
\[
  \gamma_n\bigl(\{|R_\gamma^{\mathbb{C}} v|_{\mathbb{C}^n} > t\}\bigr)
    \;\leq\; \gamma_n(\Omega)
      + \gamma_n\bigl(\{|R_\gamma^{\mathbb{C}} w|_{\mathbb{C}^n} > t\}\bigr) .
\]
\end{enumerate}

This is the reference paper's sentence ``The transformed bad part vanishes
outside $\Omega$'' in the complex setting, and it is the complex counterpart of
\noderef{TransformedLocalization}.  It is not a restatement of that node: here
there are \emph{two} bad parts $u_1, u_2$, sharing a \emph{single} exceptional
set $\Omega$, and the conclusion is about the complex operator
$R_\gamma^{\mathbb{C}}$ and the norm of $\mathbb{C}^n$, which is not a level set
of either real transform.  As in the real node, there is no splitting of the
level into $t/2 + t/2$: off $\Omega$ the two complex outputs are exactly equal.
\end{lemma}
\begin{proof}
Write $\iota : \mathbb{R}^n \to \mathbb{C}^n$ for the coordinatewise inclusion
of \noderef{ComplexRieszTransform}, so that by that node
\begin{equation*}
  R_\gamma^{\mathbb{C}} v
    = \iota\bigl(R_\gamma(\operatorname{Re} v)\bigr)
      + i\,\iota\bigl(R_\gamma(\operatorname{Im} v)\bigr),
  \qquad
  R_\gamma^{\mathbb{C}} w
    = \iota\bigl(R_\gamma(\operatorname{Re} w)\bigr)
      + i\,\iota\bigl(R_\gamma(\operatorname{Im} w)\bigr) . \tag{D}
\end{equation*}
By the hypotheses $\operatorname{Re} w = g_1$ and $\operatorname{Im} w = g_2$, the
second formula reads
$R_\gamma^{\mathbb{C}} w = \iota(R_\gamma g_1) + i\,\iota(R_\gamma g_2)$.

\emph{Step 1: the two real localizations.}
Apply \noderef{TransformedLocalization} with the data
$(f, g, u) = (\operatorname{Re} v,\, g_1,\, u_1)$, this $t$ and this $\Omega$.
Its hypotheses are met: $u_1 \in \operatorname{Dom}(L^{1/2})$ is hypothesis
$hu_1$ here, $\operatorname{Re} v = g_1 + Bu_1$ is hypothesis $h_1$, and
$Du_1 = 0$ a.e.\ on $\Omega^c$ is the first locality hypothesis.  Its first
clause gives
\begin{equation*}
  R_\gamma(\operatorname{Re} v) = R_\gamma g_1 + Du_1
  \qquad \gamma_n\text{-a.e.} \tag{1}
\end{equation*}
Applying the same node with $(f,g,u) = (\operatorname{Im} v,\, g_2,\, u_2)$, the
same $t$ and the same $\Omega$, whose hypotheses are $hu_2$, $h_2$ and the
second locality hypothesis, gives
\begin{equation*}
  R_\gamma(\operatorname{Im} v) = R_\gamma g_2 + Du_2
  \qquad \gamma_n\text{-a.e.} \tag{2}
\end{equation*}

\emph{Step 2: a single null set.}
Let $N$ be the union of the following five $\gamma_n$-null sets: the set where
$x \notin \Omega$ but $Du_1(x) \neq 0$; the set where $x \notin \Omega$ but
$Du_2(x) \neq 0$; the sets where $(1)$ and where $(2)$ fail; and the set where
the two formulas $(\mathrm{D})$ fail for the chosen representatives (each of the
operations in $(\mathrm{D})$ --- applying $\iota$ pointwise, adding, and
multiplying by $i$ --- is performed on representatives and so induces an
identity valid outside a null set).  A finite union of null sets is null
(subadditivity of a measure, shared context), so $\gamma_n(N) = 0$.

\emph{Step 3: clause (1).}
Let $x \notin \Omega \cup N$.  Then $Du_1(x) = 0$ and $Du_2(x) = 0$, so $(1)$
and $(2)$ at this $x$ give
\[
  R_\gamma(\operatorname{Re} v)(x) = R_\gamma g_1(x),
  \qquad
  R_\gamma(\operatorname{Im} v)(x) = R_\gamma g_2(x) .
\]
Substituting into the first formula of $(\mathrm{D})$ at this $x$, and then
reading the second formula backwards,
\[
  R_\gamma^{\mathbb{C}} v(x)
    = \iota\bigl(R_\gamma g_1(x)\bigr) + i\,\iota\bigl(R_\gamma g_2(x)\bigr)
    = R_\gamma^{\mathbb{C}} w(x) .
\]
Since $N$ is null, this holds for $\gamma_n$-almost every $x$ with
$x \notin \Omega$, which is clause (1).

\emph{Step 4: clause (2).}
Let $x$ satisfy $|R_\gamma^{\mathbb{C}} v(x)|_{\mathbb{C}^n} > t$ and
$x \notin \Omega \cup N$.  By Step 3,
$R_\gamma^{\mathbb{C}} v(x) = R_\gamma^{\mathbb{C}} w(x)$, so
$|R_\gamma^{\mathbb{C}} w(x)|_{\mathbb{C}^n} > t$ as well.  Hence
\[
  \{|R_\gamma^{\mathbb{C}} v|_{\mathbb{C}^n} > t\}
    \;\subseteq\; \Omega \cup N \cup
      \{|R_\gamma^{\mathbb{C}} w|_{\mathbb{C}^n} > t\} ,
\]
that is, $\{|R_\gamma^{\mathbb{C}} v| > t\}$ is contained in
$\Omega \cup \{|R_\gamma^{\mathbb{C}} w| > t\}$ up to the null set $N$.  By
monotonicity of $\gamma_n$ up to null sets and then subadditivity (both shared
context),
\[
  \gamma_n\bigl(\{|R_\gamma^{\mathbb{C}} v| > t\}\bigr)
    \leq \gamma_n\bigl(\Omega \cup \{|R_\gamma^{\mathbb{C}} w| > t\}\bigr)
    \leq \gamma_n(\Omega)
      + \gamma_n\bigl(\{|R_\gamma^{\mathbb{C}} w| > t\}\bigr) ,
\]
which is clause (2).
\end{proof}
